\documentclass[10pt]{amsart}
\usepackage[a4paper,margin=23mm]{geometry}
\usepackage{amsmath,amssymb,amsthm,mathtools,hyperref,graphicx,xspace,letltxmacro,xpatch}
\usepackage[english]{babel}
\newtheorem{theo}{Theorem}[section]
\newtheorem{lemm}[theo]{Lemma}\newtheorem{prop}[theo]{Proposition}\newtheorem{coro}[theo]{Corollary}\newtheorem{conj}[theo]{Conjecture}
\theoremstyle{definition}\newtheorem{defi}[theo]{Definition}\newtheorem{rema}[theo]{Remark}\newtheorem{exam}[theo]{Example}\newtheorem{exem}[theo]{Example}\newtheorem{nota}[theo]{Notation}
\newcommand{\ie}{i.e.}\newcommand{\eg}{e.g.}\newcommand{\resp}{resp.}\newcommand{\skpt}{\par\smallskip\noindent}
\emergencystretch=3em\pagestyle{plain}
% Add the option ``french'' if the article is in French
% Remove the option ``Unicode'' if your file is latin1-encoded


% Uncomment next line if the editor is a woman
% \womandeditor



\def\arXiv#1{\href{https://arxiv.org/abs/#1}{arXiv\string:\allowbreak#1}}


\graphicspath{{./figures/}}

\newtheorem{theorem}[theo]{Theorem}\newtheorem{thm}[theo]{Theorem}\newtheorem{lemma}[theo]{Lemma}\newtheorem{proposition}[theo]{Proposition}\newtheorem{corollary}[theo]{Corollary}\newtheorem{remark}[theo]{Remark}\newtheorem{definition}[theo]{Definition}
